\documentclass[10pt,a4paper]{article}

\usepackage[utf8]{inputenc}
\usepackage[T1]{fontenc}
\usepackage[english]{babel}
\usepackage{lmodern}
\usepackage{microtype}
\usepackage{geometry}
\usepackage{amsmath,amssymb,amsthm,mathtools,bm}
\usepackage{booktabs}
\usepackage{enumitem}
\usepackage{xcolor}
\usepackage{hyperref}
\Urlmuskip=0mu plus 1mu\relax

\geometry{margin=1.6cm}
\hypersetup{
  colorlinks=true,
  linkcolor=blue!55!black,
  citecolor=blue!55!black,
  urlcolor=blue!55!black,
  pdftitle={Non-Euclidean Pandrosion IRP},
  pdfauthor={Ivan Besevic}
}

\setlist[itemize]{leftmargin=1.8em,itemsep=0.25em,topsep=0.35em}
\setlist[enumerate]{leftmargin=1.8em,itemsep=0.25em,topsep=0.35em}
\setlength{\emergencystretch}{2em}

\theoremstyle{plain}
\newtheorem{theorem}{Theorem}[section]
\newtheorem{proposition}[theorem]{Proposition}
\newtheorem{lemma}[theorem]{Lemma}
\newtheorem{corollary}[theorem]{Corollary}

\theoremstyle{definition}
\newtheorem{definition}[theorem]{Definition}
\newtheorem{example}[theorem]{Example}

\theoremstyle{remark}
\newtheorem{remark}[theorem]{Remark}

\newcommand{\R}{\mathbb{R}}
\newcommand{\C}{\mathbb{C}}
\newcommand{\N}{\mathbb{N}}
\newcommand{\dd}{\mathrm{d}}
\newcommand{\Exp}{\operatorname{Exp}}
\newcommand{\Log}{\operatorname{Log}}
\newcommand{\dist}{\operatorname{dist}}
\newcommand{\Sp}{S_p}
\newcommand{\IRP}{\mathcal I}
\newcommand{\cB}{\mathcal B}
\newcommand{\cL}{\mathcal L}
\newcommand{\cU}{\mathcal U}
\newcommand{\Id}{\operatorname{Id}}
\newcommand{\Mob}{\operatorname{Mob}}

\newenvironment{takeaway}{%
  \begin{center}\begin{minipage}{0.92\linewidth}\small\itshape\color{blue!45!black}%
}{%
  \end{minipage}\end{center}%
}

\title{\bfseries Non-Euclidean Pandrosion IRP\\[-0.1em]
\large Geodesic Logarithms, Hadamard Charts, and Curvature Terms}
\author{Ivan \textsc{Besevic}\\[-0.2em]
\small Independent Mathematical Research\\[-0.2em]
\small \texttt{github.com/ivan-fr/poussiere\_de\_x}}
\date{June 2026}

\begin{document}
\maketitle

\begin{abstract}
The monomial Pandrosion paper shows that the raw map
\[
  h_{p,y}(s)=1-\frac{y-1}{y(1+s+\cdots+s^{p-1})}
\]
has fixed point $y^{-1/p}$ and exact multiplier
\[
  \lambda_{p,y}=1-\frac{p}{S_p(y^{1/p})}.
\]
Its acceleration principle is logarithmic: before applying the raw layer, choose
a chart and a scale that make $|\log y|$ small.  This note extracts the
non-Euclidean content of that principle.  On a general Riemannian manifold the
expression $z^p$ has no canonical meaning, so a rigorous extension must choose
extra structure.  The safest extension is the Hadamard/geodesic-dilation model:
replace the scalar logarithm by the Riemannian logarithm, replace inversion by
geodesic symmetry, and replace scalar scale cells by cells in a tangent
logarithmic chart.  In that model one renormalized raw Pandrosion layer remains
quadratic in the geodesic logarithmic residual, and $k$ layers have dyadic local
order $2^k$.  Hyperbolic space gives an explicit Poincare-ball implementation.
Lie groups and symmetric spaces give a richer, noncommutative version in which
curvature and Baker--Campbell--Hausdorff commutators enter as higher-order
defect terms.
\end{abstract}

\begin{takeaway}
\textbf{Main research translation.}
Pandrosion IRP does not need Euclidean straight lines; it needs a logarithmic
chart in which residuals are measured by a norm, and a chart-scaling rule that
keeps this residual norm small before the raw monomial layer acts.
\end{takeaway}

\section{Starting Point From The Monomial Paper}

The monomial scheme is organized by three facts:
\begin{enumerate}
  \item The proportional-mean denominator is
  \[
    S_p(s)=1+s+\cdots+s^{p-1}.
  \]
  \item The raw multiplier at the normalized target $y$ is
  \[
    \lambda_{p,y}
      =1-\frac{p}{S_p(y^{1/p})},
      \qquad y\ge 1.
  \]
  It is strictly increasing in $\log y$.
  \item Near the ideal target $y=1$,
  \[
    \lambda_{p,e^\delta}
      =\frac{p-1}{2p}\delta+O_p(\delta^2).
  \]
\end{enumerate}
Thus the geometry does not merely normalize magnitudes.  It preconditions the
multiplier itself.  IRP repeats this operation dynamically by recomputing the
residual chart after each raw layer.

\section{The Necessary Obstruction}

\begin{remark}[No canonical power on a manifold]
On a general non-Euclidean manifold $M$, there is no intrinsic operation
$z\mapsto z^p$.  A monomial problem therefore needs one of the following extra
structures:
\begin{itemize}
  \item a chosen base point and geodesic dilation, giving
  $D_p^o(z)=\Exp_o(p\Log_o z)$;
  \item a group or symmetric-space structure, giving a meaningful group power
  $z^p$ or a Cartan/geodesic power;
  \item a concrete model, such as the Poincare ball, where geodesic scalar
  multiplication has closed form.
\end{itemize}
Without one of these choices, ``Pandrosion for $z^p=x$'' is not a well-defined
mathematical statement.
\end{remark}

The rest of this note uses the first option as the rigorous baseline and records
the second option as the noncommutative research direction.

\section{Hadamard Geodesic-Dilation Model}

Let $(M,g)$ be a Hadamard manifold: complete, simply connected, with nonpositive
sectional curvature.  Then each exponential map $\Exp_o:T_oM\to M$ is globally
well behaved in the Cartan--Hadamard setting.  The positive multiplicative line
used by the original paper is the flat special case after the coordinate change
$y=e^\xi$.

\begin{definition}[Geodesic monomial]
Fix $o\in M$ and $p\ge2$.  Define the geodesic $p$-dilation
\[
  D_p^o(z)=\Exp_o(p\Log_o z).
\]
The Hadamard monomial root problem is
\[
  D_p^o(z)=x.
\]
If $X=\Log_o x$, its exact root is
\[
  r=\Exp_o(X/p).
\]
\end{definition}

This problem is deliberately narrow, but it is the cleanest non-Euclidean
analogue of the logarithmic root problem.  It replaces $\log y$ by the vector
$X=\Log_o x$.

\begin{definition}[Raw scalar Pandrosion length layer]
For a residual length $\delta\ge0$, define
\[
  s_p(\delta)
    =h_{p,e^\delta}(1)
    =1-\frac{1-e^{-\delta}}{p},
\]
and
\[
  \Theta_p(\delta)
    =\log\!\left(e^\delta s_p(\delta)^{p-1}\right)
    =\delta+(p-1)\log s_p(\delta).
\]
The number $\Theta_p(\delta)$ is the logarithmic correction produced by one raw
Pandrosion layer from the neutral point in the residual chart.
\end{definition}

\begin{lemma}[Scalar expansion]
As $\delta\to0^+$,
\[
  \Theta_p(\delta)
    =\frac{\delta}{p}
      +\frac{(p-1)^2}{2p^2}\delta^2
      +O_p(\delta^3).
\]
Consequently,
\[
  \delta-p\Theta_p(\delta)
    =-\frac{(p-1)^2}{2p}\delta^2+O_p(\delta^3).
\]
\end{lemma}

\begin{proof}
Use
\[
  s_p(\delta)=1-\frac{\delta}{p}+\frac{\delta^2}{2p}+O_p(\delta^3)
\]
and expand $\log(1+t)$.  This gives
\[
  \log s_p(\delta)
    =-\frac{\delta}{p}
      +\frac{p-1}{2p^2}\delta^2
      +O_p(\delta^3),
\]
which proves both displayed estimates.
\end{proof}

\begin{definition}[Continuous Hadamard IRP layer]
Let $u=\Exp_o(U)$ be the current root-coordinate approximation and put
\[
  A(u)=X-pU\in T_oM,\qquad \delta(u)=\|A(u)\|_g.
\]
If $\delta(u)>0$, one continuous-scale Hadamard IRP layer is
\[
  U_+
    =U+\Theta_p(\delta(u))\,\frac{A(u)}{\delta(u)},
  \qquad
  \IRP_p^o(u)=\Exp_o(U_+).
\]
If $\delta(u)=0$, set $\IRP_p^o(u)=u$.
\end{definition}

\begin{theorem}[Quadratic residual collapse in the Hadamard log chart]
For the geodesic-dilation problem $D_p^o(z)=x$, the continuous Hadamard IRP layer
satisfies
\[
  \|A(\IRP_p^o(u))\|_g
    =\frac{(p-1)^2}{2p}\,\|A(u)\|_g^2
      +O_p(\|A(u)\|_g^3).
\]
Equivalently, if the root logarithmic error is
\[
  E(u)=U-X/p,
\]
then
\[
  \|E(\IRP_p^o(u))\|_g
    =\frac{(p-1)^2}{2}\,\|E(u)\|_g^2
      +O_p(\|E(u)\|_g^3).
\]
\end{theorem}

\begin{proof}
Since $A(u)=X-pU=-pE(u)$, the update changes $A$ by
\[
  A_+
    =X-pU_+
    =A-p\Theta_p(\|A\|)\frac{A}{\|A\|}.
\]
Taking norms and applying the scalar expansion gives the result.
\end{proof}

\begin{corollary}[Dyadic local order]
For every fixed $k\ge1$, the $k$-fold local composition satisfies
\[
  \|E((\IRP_p^o)^{\circ k}(u))\|_g
    =O_{p,k}(\|E(u)\|_g^{2^k}).
\]
\end{corollary}

\begin{proof}
The previous theorem gives $\|E_+\|\le C\|E\|^2$ in a sufficiently small
neighborhood of the root.  Iteration gives the displayed dyadic order.
\end{proof}

\section{Geodesic Chart Scaling}

The continuous layer above assumes the residual length is already small.  The
non-Euclidean analogue of Thales/Riemann scaling is to reduce the residual
inside $T_oM$ before applying the raw layer.

\begin{definition}[Logarithmic scale palette in a tangent chart]
Let $\cL\subset T_oM$ be a finite or countable scale palette.  For a residual
$A=X-pU$, choose $B\in\cL$ and work with the normalized residual
\[
  A_B=A-pB.
\]
The chart-scaled layer is
\[
  U_+
    =U+B+\Theta_p(\|A_B\|)\frac{A_B}{\|A_B\|},
\]
with the evident convention when $A_B=0$.
\end{definition}

The geodesic inversion analogue is the symmetry at the chart center:
\[
  \sigma_o(\Exp_o V)=\Exp_o(-V).
\]
Thus the direct chart sees $A$, and the inverted chart sees $-A$.  Selecting the
chart and scale means minimizing $\|A_B\|$ or $\|-A-pB\|$ among admissible
choices.

\begin{proposition}[Radial two-chart lattice bound]
Fix a geodesic ray and a root-scale spacing $h=\log q>0$.  Let admissible scales
on that ray be $B_m=mh\,\nu$, $m\in\mathbb Z$, where $\nu$ is a unit tangent
direction.  For a radial residual $A=a\nu$, the best direct scale leaves a
remainder $r\in[0,ph)$, while allowing the geodesic symmetry leaves a remainder
\[
  r_*=\min(r,ph-r).
\]
Hence the worst normalized residual length is at most $ph/2$, and the exact raw
multiplier is bounded by
\[
  1-\frac{p}{S_p(e^{h/2})}.
\]
With $K$ equally spaced offsets on the same logarithmic period, the sharp
worst-case bound becomes
\[
  1-\frac{p}{S_p(e^{h/(2K)})}
  =
  1-\frac{p}{S_p(q^{1/(2K)})}.
\]
\end{proposition}

\begin{proof}
This is the same circle-covering proof as in the monomial paper, now applied to
the scalar coordinate $a$ along the geodesic ray.  The multiplier bound follows
from the monotonicity of $\lambda_{p,e^\delta}$ in $\delta$ and the identity
$\lambda_{p,e^\delta}=1-p/S_p(e^{\delta/p})$.
\end{proof}

\section{Hyperbolic Poincare-Ball Realization}

The most concrete non-Euclidean model is the Poincare ball
\[
  \mathbb B^n=\{x\in\R^n:\|x\|<1\}
\]
with curvature $-1$ and center $0$.  Its radial exponential and logarithm are
\[
  \Exp_0(V)=\tanh(\|V\|/2)\frac{V}{\|V\|},
  \qquad
  \Log_0(x)=2\,\operatorname{artanh}(\|x\|)
          \frac{x}{\|x\|}.
\]
Therefore geodesic scalar dilation is
\[
  D_a^0(x)
    =\Exp_0(a\Log_0 x)
    =
    \tanh\!\left(a\,\operatorname{artanh}\|x\|\right)\frac{x}{\|x\|}.
\]
This is the Mobius/gyrovector scalar multiplication used in hyperbolic models.

\begin{example}[Hyperbolic Pandrosion IRP layer]
Given target $x\in\mathbb B^n$ and current $u\in\mathbb B^n$:
\begin{enumerate}
  \item Compute $X=\Log_0 x$ and $U=\Log_0 u$.
  \item Set $A=X-pU$ and $\delta=\|A\|$.
  \item Optionally choose a tangent scale $B$ and replace $A$ by $A-pB$.
  \item Update
  \[
    U_+=U+B+\Theta_p(\|A-pB\|)
      \frac{A-pB}{\|A-pB\|}.
  \]
  \item Return $u_+=\Exp_0(U_+)$.
\end{enumerate}
This solves the hyperbolic geodesic root problem
$D_p^0(u)=x$, not the Euclidean coordinate equation $u^p=x$.
\end{example}

\section{Beyond The Radial Model: Curvature And Noncommutativity}

The Hadamard construction above is rigorous because it fixes a logarithmic chart.
It becomes genuinely curved when the chart center moves, or when the monomial
power is a group/symmetric-space operation rather than a radial geodesic
dilation.

\begin{definition}[Lie-group residual IRP, local form]
Let $G$ be a Lie group with exponential map $\exp:\mathfrak g\to G$.  For the
local group-power problem $z^p=x$, define a right residual
\[
  R(u)=u^{-p}x,\qquad A(u)=\log R(u)\in\mathfrak g.
\]
A residual Pandrosion layer applies the scalar length correction in the residual
direction:
\[
  u_+=u\,\exp\!\left(
    \Theta_p(\|A(u)\|)
    \frac{A(u)}{\|A(u)\|}
  \right),
\]
with scale and inversion choices made in $\mathfrak g$.
\end{definition}

\begin{remark}[Where curvature appears]
In the abelian case, the preceding formula reduces to the scalar logarithmic IRP.
In the noncommutative case,
\[
  \log(\exp A\,\exp B)
    =A+B+\frac12[A,B]+\cdots
\]
introduces Baker--Campbell--Hausdorff commutators.  On symmetric spaces the same
role is played by curvature/Jacobi-field terms.  Thus the local research target
is not the scalar quadratic term, which is already known, but the estimate
\[
\begin{aligned}
  \|A_+\|
    &\le C_p\|A\|^2
      +C_{\mathrm{curv}}\|A\|\,\|[A,U]\|  \\
    &\quad +O(\|A\|^3).
\end{aligned}
\]
If the residual direction nearly commutes with the current chart, or if the
problem is radial in a symmetric space, the commutator term vanishes and the
Hadamard theorem above is recovered.
\end{remark}

\begin{proposition}[Local curvature-stable template]
Let a moving-chart non-Euclidean IRP layer be written in normal coordinates at a
root $r$ as
\[
  E_+=Q_p(E)+C_\kappa(E),
\]
where $Q_p(E)$ is the flat Pandrosion residual map with
$\|Q_p(E)\|\le c_p\|E\|^2$ and the curvature defect satisfies
$\|C_\kappa(E)\|\le c_\kappa\|E\|^3$ on a normal ball.  Then the layer remains
quadratic:
\[
  \|E_+\|\le C\|E\|^2.
\]
Consequently fixed $k$-fold IRP compositions still have local order $2^k$.
\end{proposition}

\begin{proof}
For $\|E\|$ sufficiently small, the cubic defect is dominated by the quadratic
term.  The dyadic-order argument is the same as in the flat and Hadamard
settings.
\end{proof}

\section{Research Program}

\begin{enumerate}
  \item \textbf{Formal baseline.}  Prove the Hadamard geodesic-dilation theorem
  in Lean or another proof assistant.  This requires only normal coordinates,
  vector norms, and the scalar expansion of $\Theta_p$.
  \item \textbf{Hyperbolic implementation.}  Implement the Poincare-ball layer
  with $\Exp_0$, $\Log_0$, and tangent scale palettes.  Verify exact reduction
  to the monomial paper on every geodesic ray.
  \item \textbf{Spherical local model.}  Repeat the construction on the sphere
  only inside injectivity-radius balls.  The cut locus replaces the global
  Hadamard guarantee, so atlas switching is mandatory.
  \item \textbf{Symmetric positive-definite model.}  Use the affine-invariant
  SPD geometry, where logarithms and geodesics are explicit.  Radial spectral
  problems are exact; noncommuting updates measure the first serious curvature
  defect.
  \item \textbf{Lie-group BCH bounds.}  Quantify when the noncommutative group
  residual keeps the same quadratic IRP order.  The key estimate is the size of
  the commutator defect relative to the scalar Pandrosion residual collapse.
\end{enumerate}

\section{Conclusion}

The non-Euclidean version of Pandrosion IRP should not be phrased as a naive
replacement of straight lines by curved lines.  The correct invariant object is
the logarithmic residual.  In the original monomial paper this residual is the
number $\log y$ on the positive multiplicative line.  In a Hadamard manifold it
becomes $\Log_o x-p\Log_o u$.  Chart scaling means minimizing the norm of that
residual before the raw layer acts.  Under this translation, the local IRP
mechanism survives unchanged: one renormalized raw layer is quadratic, and
iterated layers have dyadic order.  Curvature enters only when the chart moves
or when the underlying power law is noncommutative; those terms form the next
research problem.

\begin{thebibliography}{9}

\bibitem{besevic000}
I. Besevic,
\emph{A Monomial Pandrosion Scheme: Contraction, Chart Scaling, and Iterated
Renormalization for $z^p=x$}, 2026.

\bibitem{bridsonhaefliger}
M. R. Bridson and A. Haefliger,
\emph{Metric Spaces of Non-Positive Curvature},
Springer, 1999.
\url{https://doi.org/10.1007/978-3-662-12494-9}

\bibitem{bacak}
M. Bacak,
\emph{Convex Analysis and Optimization in Hadamard Spaces},
De Gruyter, 2014.
\url{https://doi.org/10.1515/9783110361629}

\bibitem{absil}
P.-A. Absil, R. Mahony, and R. Sepulchre,
\emph{Optimization Algorithms on Matrix Manifolds},
Princeton University Press, 2008.
\url{https://sites.uclouvain.be/absil/amsbook/}

\bibitem{nickelkiela}
M. Nickel and D. Kiela,
\emph{Poincare Embeddings for Learning Hierarchical Representations},
Advances in Neural Information Processing Systems 30, 2017.
\url{https://papers.neurips.cc/paper/7213-poincare-embeddings-for-learning-hierarchical-representations}

\end{thebibliography}

\end{document}
